\chapter{Geometri Pemetaan Linear}

Untuk memperlihatkan efek geometris transformasi linear,
kita akan menggambarkan aksi
transformasi bidang pada persegi satuan.
Kita menampilkan bidang, $\Re^2$,
karena dapat dimuat di atas kertas, tetapi prinsip-prinsipnya berlaku pula pada
ruang berdimensi lebih tinggi~$\Re^n$.




%========================================
\section{Garis dipetakan ke garis}
Bab sebelumnya menunjukkan bahwa
syarat $h(r\cdot \vec{v})=r\cdot h(\vec{v})$ dalam definisi
pemetaan linear berarti bahwa
pemetaan semacam ini membawa garis yang melalui titik asal ke garis yang melalui titik asal.
Bagaimana dengan garis yang tidak melalui titik asal?
Tetapkan suatu pemetaan linear $\map{h}{\Re^d}{\Re^c}$.
Suatu garis dalam domain berbentuk
$\ell = \set{\vec{m}\cdot s +\vec{b}\suchthat s\in\Re}$
untuk~$\vec{m}, \vec{b}\in\Re^d$ (jika $\vec{m}=\zero$, maka garis itu degenerat).
Citranya
\begin{equation*}
  h(\ell)=\set{h(\vec{m}\cdot s+\vec{b})\suchthat s\in\Re}
  =\set{h(\vec{m})\cdot s+h(\vec{b})\suchthat s\in\Re}
\end{equation*}
adalah suatu garis dalam kodomain.
Jadi, pemetaan linear membawa setiap garis, baik melalui titik asal maupun tidak, ke garis lain.

Sebagai contoh, perhatikan transformasi $\map{t}{\Re^2}{\Re^2}$
yang memutar vektor berlawanan arah jarum jam sebesar
$\pi/6$~radian.
\begin{equation*}
  \rep{t}{\stdbasis_2,\stdbasis_2}
  =
  \begin{mat}
    \cos(\pi/6)  &-\sin(\pi/6) \\
    \sin(\pi/6)  &\cos(\pi/6)
  \end{mat}
  = 
  \begin{mat}
    \sqrt{3}/2   &-1/2 \\
    1/2          &\sqrt{3}/2
  \end{mat}
\end{equation*}
Perhatikan pula garis $y=3x+2$.
\begin{equation*}
  \ell=
  \set{\colvec{x \\ y}=\colvec{1 \\ 3}\cdot s+\colvec{0 \\ 2}\suchthat s\in\Re}
\end{equation*}
Di bawah aksi~$t$,
\begin{equation*}
  \begin{mat}
    \sqrt{3}/2   &-1/2 \\
    1/2          &\sqrt{3}/2
  \end{mat}
  \colvec{1  \\ 3}
  =
  \colvec{(\sqrt{3}-3)/2  \\ (1+3\sqrt{3})/2}
  \qquad
  \begin{mat}
    \sqrt{3}/2    &-1/2 \\
    1/2          &\sqrt{3}/2
  \end{mat}
  \colvec{0  \\ 2}
  =
  \colvec{-1  \\ \sqrt{3}}
\end{equation*}
garis tersebut menjadi himpunan berikut.
\begin{equation*}
  t(\ell)
  =
  \set{\colvec{x \\ y}=\colvec{(\sqrt{3}-3)/2 \\ (1+3\sqrt{3})/2}\cdot s
                          +\colvec{-1 \\ \sqrt{3}}\suchthat s\in\Re}
\end{equation*}
Berikut adalah kode serta plot masukan dan keluarannya.\footnote{%
  Seperti dalam bab-bab sebelumnya, sebagian kode yang
  mengatur ukuran huruf dan sebagainya tidak ditampilkan.}
  %% For the full call see this manual's source.
  %% Also, the limits on the parameter~$s$ of $10$ and~$-10$ are arbitrary, just
  %% chosen to be large enough that the line segment covers the entire 
  %% domain and codomain intervals shown, from $-4$ to~$4$.}
\begin{sagecommandline}
sage: s = var('s')
sage: ell = parametric_plot((s, 3*s+2), (s, -10, 10))
sage: d = circle((0,2), 0.1, rgbcolor=(1,0,0))
sage: ell = ell+d
sage: ell.save("graphics/geo000a.pdf")
sage: t_x(s) = ((sqrt(3)-3)/2)*s-1
sage: t_y(s) = ((1+3*sqrt(3))/2)*s+sqrt(3)
sage: t_ell = parametric_plot((t_x(s), t_y(s)), (s, -10, 10))
sage: d = circle((-1,sqrt(3)), 0.1, rgbcolor=(1,0,0))
sage: t_ell = t_ell+d
sage: t_ell.save("graphics/geo000b.pdf")
\end{sagecommandline}
\begin{sagesilent}
s = var('s')
plot.options['figsize'] = 2.5
plot.options['axes_pad'] = 0.05
plot.options['fontsize'] = 7
plot.options['dpi'] = 1200
plot.options['aspect_ratio'] = 1
ell = parametric_plot((s, 3*s+2), (s, -10, 10))
d = circle((0,2), 0.1, rgbcolor=(1,0,0))
ell = ell+d
ell.set_axes_range(-4, 4, -4, 4)
ell.save("graphics/geo000a.pdf", fontsize=7)
t_x(s) = ((sqrt(3)-3)/2)*s-1
t_y(s) = ((1+3*sqrt(3))/2)*s+sqrt(3)
t_ell = parametric_plot((t_x(s), t_y(s)), (s, -10, 10))
d = circle((-1,sqrt(3)), 0.1, rgbcolor=(1,0,0))
t_ell = t_ell+d
t_ell.set_axes_range(-4, 4, -4, 4)
t_ell.save("graphics/geo000b.pdf", fontsize=7)
\end{sagesilent}
\begin{equation*}
  \vcenteredhbox{\includegraphics{graphics/geo000a.pdf}}
  \quad\mapsvia{\makebox[1.5em][c]{\scriptsize $t$}}\quad
  \vcenteredhbox{\includegraphics{graphics/geo000b.pdf}}
\end{equation*}
Pemetaan~$t$ merupakan rotasi sebesar~$\pi/6$~radian, tetapi hal itu mungkin sulit terlihat.
Salah satu contohnya, yang ditunjukkan oleh titik-titik merah, ialah bahwa vektor yang berujung di $(0,2)$
diputar menjadi vektor yang berujung di $(-1,\sqrt{3})$.




%========================================
\section{Persegi satuan}
Bab sebelumnya menggambarkan
efek transformasi bidang
$\map{t}{\Re^2}{\Re^2}$ 
dengan menerapkannya pada lingkaran satuan.
Dalam bab ini kita menggunakan persegi satuan.
Pengamatan bahwa
pemetaan linear membawa garis ke garis membuat kodenya sederhana:~kita mulai dengan
keempat sudut persegi masukan,
menghitung efek transformasi pada sudut-sudut itu untuk
mendapatkan empat sudut keluaran, lalu
menghubungkan keempatnya dengan ruas garis sehingga diperoleh jajaran genjang.
(Persamaan-persamaan di atas menunjukkan bahwa garis-garis
dengan vektor arah~$\vec{m}$ yang sama dipetakan ke garis-garis dengan vektor arah
$t(\vec{m})$ yang sama; dengan kata lain, garis-garis sejajar dipetakan ke garis-garis sejajar.)

Jadi, metode kita adalah memulai dengan persegi satuan.
Perhatikan warnanya; di sini warna-warna itu berperan sama seperti sebelumnya.
\begin{sagecommandline}
sage: load("plot_action.sage")
sage: p = plot_square_action(1,0,0,1)  # matriks identitas
sage: p.save("graphics/geo100.pdf")
\end{sagecommandline}
\begin{sagesilent}
load("plot_action.sage")
p = plot_square_action(1,0,0,1)  # matriks identitas
p.set_axes_range(-0.25, 1.5, -0.25, 1.5) 
p.save("graphics/geo100.pdf")
\end{sagesilent}
\begin{center}
  \includegraphics{graphics/geo100.pdf}
\end{center}
Rutin \inlinecode{plot_square_action(a,b,c,d)} menerapkan pada
persegi satuan tersebut suatu
transformasi yang, terhadap basis standar, direpresentasikan oleh
matriks dengan entri $a$, $b$, $c$, dan~$d$.
(Kode sumbernya terdapat pada akhir bab.)
\begin{equation*}
  \begin{mat}
    a &b \\
    c &d
  \end{mat}
  \colvec{x  \\ y}
  =
  \colvec{ax+by \\ cx+dy}
\end{equation*}
Berikut ini menggambarkan efek matriks umum tersebut.
\begin{sagecommandline}
sage: load("plot_action.sage")
sage: q = plot_square_action(1,0,0,1) 
sage: q.set_axes_range(-1, 4, -1, 7) 
sage: q.save("graphics/geo101a.pdf")
sage: p = plot_square_action(1,2,3,4) 
sage: p.set_axes_range(-1, 4, -1, 7) 
sage: p.save("graphics/geo101b.pdf")
\end{sagecommandline}
% \begin{sagesilent}
% load("plot_action.sage")
% q = plot_square_action(1,0,0,1) 
% q.set_axes_range(-1, 4, -1, 7) 
% q.save("graphics/geo101a.pdf")
% p = plot_square_action(1,2,3,4) 
% p.set_axes_range(-1, 4, -1, 7) 
% p.save("graphics/geo101b.pdf")
% \end{sagesilent}
\begin{equation*}
  \vcenteredhbox{\includegraphics{graphics/geo101a.pdf}}
  \quad\mapsvia{\big (\begin{smallmatrix} 1 &2 \\ 3 &4 \end{smallmatrix}\big )}\quad
  \vcenteredhbox{\includegraphics{graphics/geo101b.pdf}}
\end{equation*}
Seperti dijelaskan di atas, transformasi itu mengubah persegi masukan menjadi jajaran genjang keluaran.

Warna-warna tersebut memperlihatkan efek lain transformasi, selain perubahan bentuk.
Jika warna diurutkan secara alami sebagai merah, jingga,
hijau, lalu biru,
persegi dalam domain berorientasi berlawanan arah jarum jam.
Namun, bangun dalam kodomain berorientasi searah jarum jam.
Jadi, warna-warna itu menunjukkan bahwa transformasi ini
membalik orientasi (yang kadang disebut `arah').

Kita dapat memahami perilaku yang rumit dengan
membangunnya dari pemahaman atas perilaku yang sederhana.
Karena itu, kita mulai dengan transformasi contoh yang sama seperti dalam bab sebelumnya,
yaitu transformasi yang menggandakan komponen pertama.
%% \footnote{%
%%   As in other chapters, some of the plots are drawn using 
%%   some fiddly graphics options that are not shown.
%%   In this case, the \protect\inlinecode{p.save(...)} has
%%   the \protect\inlinecode{ticks\_integer=True} option.}
\begin{sagecommandline}
sage: load("plot_action.sage")
sage: q = plot_square_action(1,0,0,1) 
sage: q.set_axes_range(-1, 3, -1, 2) 
sage: q.save("graphics/geo102a.pdf")
sage: p = plot_square_action(2,0,0,1)  
sage: p.set_axes_range(-1, 3, -1, 2) 
sage: p.save("graphics/geo102b.pdf")
\end{sagecommandline}
\begin{sagesilent}
load("plot_action.sage")
q = plot_square_action(1,0,0,1) 
q.set_axes_range(-1, 3, -1, 2) 
q.save("graphics/geo102a.pdf", ticks_integer=True)
p = plot_square_action(2,0,0,1)  
p.set_axes_range(-1, 3, -1, 2) 
p.save("graphics/geo102b.pdf", ticks_integer=True)
\end{sagesilent}
\begin{equation*}
  \vcenteredhbox{\includegraphics{graphics/geo102a.pdf}}
  \quad\mapsvia{\big (\begin{smallmatrix} 2 &0 \\ 0 &1 \end{smallmatrix}\big )}\quad
  \vcenteredhbox{\includegraphics{graphics/geo102b.pdf}}
  \tag{$*$}
\end{equation*}
Pemetaan linear membawa vektor nol ke vektor nol, dan persegi
masukan berpangkal di titik asal, sehingga
bangun keluaran juga berpangkal di titik asal.
Namun, bangun itu telah diregangkan secara horizontal\Dash orientasinya sama
dengan persegi awal, tetapi luasnya dua kali lipat.

Contoh tersebut menunjukkan bahwa perilaku yang berkaitan dengan matriks diagonal
bersifat sederhana.
Sebagai contoh, mengalikan koordinat~$x$ dengan tiga menghasilkan bangun serupa yang
luasnya tiga kali luas bangun awal.
Apa yang terjadi jika kita mengambil $-1$ kali koordinat~$x$?
\begin{sagecommandline}
sage: load("plot_action.sage")
sage: q = plot_square_action(1,0,0,1) 
sage: q.set_axes_range(-2, 2, -1, 2) 
sage: q.save("graphics/geo103a.pdf")
sage: p = plot_square_action(-1,0,0,1) 
sage: p.set_axes_range(-2, 2, -1, 2) 
sage: p.save("graphics/geo103b.pdf")
\end{sagecommandline}
\begin{sagesilent}
load("plot_action.sage")
q = plot_square_action(1,0,0,1) 
q.set_axes_range(-2, 2, -1, 2) 
q.save("graphics/geo103a.pdf", ticks_integer=True)
p = plot_square_action(-1,0,0,1) 
p.set_axes_range(-2, 2, -1, 2) 
p.save("graphics/geo103b.pdf", ticks_integer=True)
\end{sagesilent}
\begin{equation*}
  \vcenteredhbox{\includegraphics{graphics/geo103a.pdf}}
  \quad\mapsvia{\big (\begin{smallmatrix} -1 &0 \\ 0 &1 \end{smallmatrix}\big )}\quad
  \vcenteredhbox{\includegraphics{graphics/geo103b.pdf}}
  \tag{$**$}
\end{equation*}
Transformasi itu mengubah orientasi.

Kita mengatakan bahwa bangun di sebelah kiri di atas memiliki \textit{luas berarah}
sebesar $-1$.
Alasan luas diberi tanda
adalah kekontinuan:~bayangkan kita mulai dengan bangun sebelah kanan
dari contoh sebelumnya, yaitu diagram berlabel~($*$),
lalu menggeser sisi jingga dari kanan, dari $2$ ke~$1$, ke~$0$, dan
kemudian ke~$-1$.
Luasnya turun dari~$2$ ke~$1$, lalu ke~$0$, sehingga secara alami
kita memberi bangun di atas ukuran luas~$-1$.
Kata `berarah' hanya digunakan untuk membedakan ukuran ini dari
pengertian luas yang dipelajari di sekolah dasar.
Luas dalam pengertian sekolah dasar adalah nilai mutlak dari luas berarah.

Transformasi berikutnya menggabungkan aksi pada dua sumbu,
yaitu mengalikan komponen~$y$ dengan tiga dan
komponen~$x$ dengan $-1$.
\begin{sagecommandline}
sage: load("plot_action.sage")
sage: q = plot_square_action(1,0,0,1) 
sage: q.set_axes_range(-2, 2, -1, 4) 
sage: q.save("graphics/geo104a.pdf")
sage: p = plot_square_action(-1,0,0,3) 
sage: p.set_axes_range(-2, 2, -1, 4) 
sage: p.save("graphics/geo104b.pdf")
\end{sagecommandline}
\begin{sagesilent}
load("plot_action.sage")
q = plot_square_action(1,0,0,1) 
q.set_axes_range(-2, 2, -1, 4) 
q.save("graphics/geo104a.pdf", ticks_integer=True)
p = plot_square_action(-1,0,0,3) 
p.set_axes_range(-2, 2, -1, 4) 
p.save("graphics/geo104b.pdf", ticks_integer=True)
\end{sagesilent}
\begin{equation*}
  \vcenteredhbox{\includegraphics{graphics/geo104a.pdf}}
  \quad\mapsvia{\big (\begin{smallmatrix} -1 &0 \\ 0 &3 \end{smallmatrix}\big )}\quad
  \vcenteredhbox{\includegraphics{graphics/geo104b.pdf}}
\end{equation*}
Warna-warnanya menunjukkan bahwa transformasi ini juga mengubah
orientasi, sehingga bangun baru memiliki luas berarah~$-3$.

Apa yang terjadi jika orientasi diubah dua kali?
\begin{sagecommandline}
sage: load("plot_action.sage")
sage: q = plot_square_action(1,0,0,1) 
sage: q.set_axes_range(-3, 2, -4, 2) 
sage: q.save("graphics/geo105a.pdf")
sage: p = plot_square_action(-2,0,0,-3) 
sage: p.set_axes_range(-3, 2, -4, 2) 
sage: p.save("graphics/geo105b.pdf")
\end{sagecommandline}
\begin{sagesilent}
load("plot_action.sage")
q = plot_square_action(1,0,0,1) 
q.set_axes_range(-3, 2, -4, 2) 
q.save("graphics/geo105a.pdf", ticks_integer=True)
p = plot_square_action(-2,0,0,-3) 
p.set_axes_range(-3, 2, -4, 2) 
p.save("graphics/geo105b.pdf", ticks_integer=True)
\end{sagesilent}
Plot sebelum dan sesudah
\begin{equation*}
  \vcenteredhbox{\includegraphics{graphics/geo105a.pdf}}
  \quad\mapsvia{\big (\begin{smallmatrix} -2 &0 \\ 0 &-3 \end{smallmatrix}\big )}\quad
  \vcenteredhbox{\includegraphics{graphics/geo105b.pdf}}
\end{equation*}
menunjukkan bahwa pada bangun sebelah kanan,
warna-warnanya muncul dalam urutan yang sama seperti pada orientasi
berlawanan arah jarum jam bangun asli:
merah, jingga, hijau, lalu biru.
Dengan kata lain, membalik orientasi dua kali mengembalikan
orientasi semula.
Bangun baru memiliki luas berarah~$6$.

Selanjutnya, perhatikan efek penempatan entri di luar diagonal pada matriks.
\begin{sagecommandline}
sage: load("plot_action.sage")
sage: q = plot_square_action(1,0,0,1) 
sage: q.set_axes_range(-1, 3, -1, 2) 
sage: q.save("graphics/geo106a.pdf")
sage: p = plot_square_action(1,2,0,1) 
sage: p.set_axes_range(-1, 3, -1, 2) 
sage: p.save("graphics/geo106b.pdf")
\end{sagecommandline}
\begin{sagesilent}
load("plot_action.sage")
q = plot_square_action(1,0,0,1) 
q.set_axes_range(-1, 3, -1, 2) 
q.save("graphics/geo106a.pdf", ticks_integer=True)
p = plot_square_action(1,2,0,1) 
p.set_axes_range(-1, 3, -1, 2) 
p.save("graphics/geo106b.pdf", ticks_integer=True)
\end{sagesilent}
\begin{equation*}
  \vcenteredhbox{\includegraphics{graphics/geo106a.pdf}}
  \quad\mapsvia{\big (\begin{smallmatrix} 1 &2 \\ 0 &1 \end{smallmatrix}\big )}\quad
  \vcenteredhbox{\includegraphics{graphics/geo106b.pdf}}
\end{equation*}
Transformasi ini merupakan suatu \textit{geseran}.
Sisi-sisi keluaran tidak saling tegak lurus, meskipun pasangan sisi yang berhadapan tetap
sejajar.
Aksi
\begin{equation*}
  \begin{mat}
    1  &2  \\
    0  &1
  \end{mat}
  \colvec{x \\ y}
  = \colvec{x+2y \\ y}
\end{equation*}
berarti bahwa
vektor awal dengan komponen~$y$ sebesar~$1$ digeser ke kanan sejauh~$2$, sedangkan
vektor awal dengan komponen~$y$ sebesar~$2$ digeser ke kanan sejauh~$4$, dan seterusnya.
Dengan kata lain,
vektor digeser bergantung pada seberapa jauh letaknya di atas atau di bawah
sumbu~$x$.
Bangun keluaran memiliki alas~$1$
dan tinggi~$1$, sementara orientasinya dipertahankan,
sehingga luas berarahnya~$1$.

Menempatkan nilai tak nol pada entri luar diagonal yang lain,
yaitu entri kiri bawah, memberikan efek yang sama, kecuali bahwa geserannya sejajar
dengan sumbu~$y$.
\begin{sagecommandline}
sage: load("plot_action.sage")
sage: q = plot_square_action(1,0,0,1) 
sage: q.set_axes_range(-1, 3, -3, 2) 
sage: q.save("graphics/geo107a.pdf")
sage: p = plot_square_action(1,0,-2,1) 
sage: p.set_axes_range(-1, 3, -3, 2) 
sage: p.save("graphics/geo107b.pdf")
\end{sagecommandline}
\begin{sagesilent}
load("plot_action.sage")
q = plot_square_action(1,0,0,1) 
q.set_axes_range(-1, 3, -3, 2) 
q.save("graphics/geo107a.pdf", ticks_integer=True)
p = plot_square_action(1,0,-2,1) 
p.set_axes_range(-1, 3, -3, 2) 
p.save("graphics/geo107b.pdf", ticks_integer=True)
\end{sagesilent}
\begin{equation*}
  \vcenteredhbox{\includegraphics{graphics/geo107a.pdf}}
  \quad\mapsvia{\big (\begin{smallmatrix} 1 &0 \\ -2 &1 \end{smallmatrix}\big )}\quad
  \vcenteredhbox{\includegraphics{graphics/geo107b.pdf}}
\end{equation*}
Aksi di atas
\begin{equation*}
  \colvec{x \\ y} \mapsto \colvec{x \\ -2x+y}
\end{equation*}
berarti bahwa vektor digeser bergantung pada seberapa jauh letaknya dari
sumbu~$y$.
Sebagai contoh, vektor masukan dengan komponen~$x$ sebesar~$1$ digeser sebesar~$-2$,
sedangkan jika komponen~$x$-nya~$2$, vektor itu digeser sebesar~$-4$.
Luas berarah bangun keluaran adalah~$1$.





\section{Determinan}
Buku ini menafsirkan secara geometris
syarat-syarat dalam definisi
fungsi determinan.
Buku ini menunjukkan bahwa, ketika beralih dari gambar sebelum ke gambar sesudah,
matriks-matriks transformasi tersebut
mengubah luas berarah daerah masukan dengan faktor sebesar
determinan matriks.
Pada diagram~($*$) di atas, matriksnya berdeterminan~$2$ dan menggandakan luas berarah.
Pada diagram~($**$), matriksnya mengalikan luas berarah dengan~$-1$.

Salah satu keuntungan memulai gambar sebelum/sesudah ini dengan persegi
satuan ialah bahwa, karena masukan berluas~$1$,
luas berarah bangun keluaran sama dengan
determinan matriks.

Sebagai contoh, matriks
\begin{equation*}
\begin{mat}
  1 &2 \\
  2 &4
\end{mat}
\end{equation*}
bersifat singular sehingga determinannya~$0$.
\begin{sagecommandline}
sage: load("plot_action.sage")
sage: q = plot_square_action(1,0,0,1) 
sage: q.set_axes_range(-1, 4, -1, 6) 
sage: q.save("graphics/geo200a.pdf")
sage: p = plot_square_action(1,2,2,4) 
sage: p.set_axes_range(-1, 4, -1, 6) 
sage: p.save("graphics/geo200b.pdf")
\end{sagecommandline}
\begin{sagesilent}
load("plot_action.sage")
q = plot_square_action(1,0,0,1) 
q.set_axes_range(-1, 4, -1, 6) 
q.save("graphics/geo200a.pdf")
p = plot_square_action(1,2,2,4) 
p.set_axes_range(-1, 4, -1, 6) 
p.save("graphics/geo200b.pdf")
\end{sagesilent}
\begin{equation*}
  \vcenteredhbox{\includegraphics{graphics/geo200a.pdf}}
  \quad\mapsvia{\big (\begin{smallmatrix} 1 &2 \\ 2 &4 \end{smallmatrix}\big )}\quad
  \vcenteredhbox{\includegraphics{graphics/geo200b.pdf}}
\end{equation*}
Bangun keluaran memiliki luas nol.






\section{Faktorisasi Turing, PA=LDU}
Sekarang kita akan melihat bagaimana aksi sebarang matriks dapat diuraikan menjadi
aksi-aksi yang ditunjukkan di atas.
Dengan demikian, kita memperoleh deskripsi geometris lengkap bagi sebarang pemetaan linear;
artinya, efek sebarang transformasi dapat dipahami dengan
menguraikannya menjadi serangkaian langkah sederhana.

Ingat bahwa kita dapat melakukan operasi baris dalam Metode Gauss dengan
perkalian matriks.
Sebagai contoh, perkalian dari kiri dengan matriks berikut memberikan efek operasi baris
$2\rho_1+\rho_2$.
\begin{equation*}
  \begin{mat}
    1 &0 &0 \\
    2 &1 &0 \\
    0 &0 &1
  \end{mat}
  \begin{mat}
    3 &1 &4 \\
   -6 &1 &-8 \\
    0 &-3 &2
  \end{mat}
  =
  \begin{mat}
    3 &1  &4 \\ 
    0 &3 &0 \\
    0 &-3  &2
  \end{mat}
\end{equation*}

Seperti dijelaskan dalam buku,
matriks reduksi elementer
terdiri atas tiga jenis.
Untuk suatu matriks~$H$, lakukan penskalaan baris~\( r\rho_i \)
dengan \( M_i(r)\,H \),
pertukarkan baris \( \rho_i\leftrightarrow\rho_j \) dengan \( P_{i,j}\,H \),
dan tambahkan kelipatan suatu baris ke baris lain,
\( k\rho_i+\rho_j \), dengan \( C_{i,j}(k)\,H \).
Ketiga jenis ini
diperoleh dengan menerapkan suatu operasi baris pada matriks identitas:
\begin{center}
$I\grstep{r\rho_i}M_i(r)$
\qquad
\( I\grstep{\rho_i\leftrightarrow\rho_j}P_{i,j} \)
\qquad
\( I\grstep{k\rho_i+\rho_j}C_{i,j}(k) \)
\end{center}
(dengan $r\neq 0$ dan $i\neq j$).
Sebagai contoh, paragraf sebelumnya menggunakan matriks $\nbyn{3}$
$C_{1,2}(2)$.
Kita akan berfokus pada transformasi, sehingga semua matriks ini
dianggap berukuran sama dan berbentuk persegi.

Untuk melanjutkan reduksi dengan Metode Gauss yang dimulai dalam persamaan matriks di atas,
gunakan $C_{2,3}(1)$ untuk melakukan $\rho_2+\rho_3$,
sehingga diperoleh bentuk eselon.
\begin{equation*}
  \begin{mat}
    1 &0  &0 \\
    0 &1  &0 \\
    0 &1 &1
  \end{mat}
  \begin{mat}
    1 &0 &0 \\
    2 &1 &0 \\
    0 &0 &1
  \end{mat}\cdot
  \begin{mat}
    3 &1 &4 \\
   -6 &1 &-8 \\
    0 &-3 &2
  \end{mat}
  =
  \begin{mat}
    3 &1  &4 \\ 
    0 &3  &0 \\
    0 &0  &2
  \end{mat}
  % \tag{$*$}
\end{equation*}

Perhatikan bahwa jika matriks awal sedemikian rupa sehingga
tidak diperlukan pertukaran baris,
kita cukup menggunakan operasi~\( k\rho_i+\rho_j \) dengan
$j>i$.
Matriks elementer
yang melakukan operasi tersebut disebut \definend{segitiga bawah}
karena semua entri di atas diagonal utamanya bernilai nol.
(Matriks yang semua entrinya di bawah diagonal utama bernilai nol disebut
\definend{segitiga atas}.)
% Thus the factorization so far is into a product of lower triangular elementary
% matrices and an echelon form matrix.

Kita dapat melangkah lebih jauh.
Selanjutnya, gunakan matriks diagonal untuk mengubah
entri utama pada baris-baris tak nol dari matriks bentuk eselon menjadi~$1$.
\begin{equation*}
  \begin{mat}
    1/3 &0   &0 \\
    0   &1/3 &0 \\
    0   &0   &1/2  
  \end{mat}
  \begin{mat}
    1 &0  &0 \\
    0 &1  &0 \\
    0 &1 &1
  \end{mat}
  \begin{mat}
    1 &0 &0 \\
    2 &1 &0 \\
    0 &0 &1
  \end{mat}\cdot
  \begin{mat}
    3 &1 &4 \\
   -6 &1 &-8 \\
    0 &-3 &2
  \end{mat}
  =
  \begin{mat}
    1 &1/3  &4/3 \\ 
    0 &1  &0 \\
    0 &0  &1
  \end{mat}
  \tag{$*$}
\end{equation*}
%% sage: M = matrix(QQ, [[3,1,4], [-6,1,-8], [0,-3,2]])                            
%% sage: N1 = matrix(QQ, [[1,0,0], [2,1,0], [0,0,1]])                              
%% sage: N2 = matrix(QQ, [[1,0,0], [0,1,0], [0,1,1]])                              
%% sage: N3 = matrix(QQ, [[1/3,0,0], [0,1/3,0], [0,0,1/2]])                        
%% sage: N3*N2*N1*M                                                                
%% [  1 1/3 4/3]
%% [  0   1   0]
%% [  0   0   1]
Terakhir, gunakan operasi kolom
untuk mereduksi matriks tersebut sepenuhnya menjadi
matriks identitas parsial berbentuk blok.
Berikut adalah perkalian dari kanan pada ruas kanan~($*$) untuk
menambahkan $-1/3$ kali kolom pertama ke kolom kedua.
\begin{equation*}
  \begin{mat}
    1 &1/3  &4/3 \\ 
    0 &1  &0 \\
    0 &0  &1
  \end{mat}\cdot
  \begin{mat}
    1  &-1/3  &0  \\
    0  &1     &0  \\
    0  &0     &1
  \end{mat}
  =
  \begin{mat}
    1 &0  &4/3 \\ 
    0 &1  &0 \\
    0 &0  &1
  \end{mat}
\end{equation*}
Terakhir, penambahan $-4/3$~kali kolom pertama ke kolom ketiga
menghasilkan matriks identitas.
\begin{equation*}
  \begin{mat}
    1 &1/3  &4/3 \\ 
    0 &1  &0 \\
    0 &0  &1
  \end{mat}\cdot
  \begin{mat}
    1  &-1/3  &0  \\
    0  &1     &0  \\
    0  &0     &1
  \end{mat}
  \begin{mat}
    1  &0     &-4/3  \\
    0  &1     &0  \\
    0  &0     &1
  \end{mat}
  =
  \begin{mat}
    1 &0  &0 \\ 
    0 &1  &0 \\
    0 &0  &1
  \end{mat}
   \tag{$**$}
\end{equation*}
%% sage: P1 = matrix(QQ, [[1,-1/3,0], [0,1,0], [0,0,1]])                           
%% sage: P2 = matrix(QQ, [[1,0,-4/3], [0,1,0], [0,0,1]])                           
%% sage: N3*N2*N1*M*P1*P2                                                          
%% [1 0 0]
%% [0 1 0]
%% [0 0 1]

Ringkasnya,
jika kita mulai dengan matriks $A$ yang tidak memerlukan pertukaran baris,
kita memperoleh persamaan matriks berikut.
\begin{equation*}
  D_0\cdot L_s\cdots L_2 L_1\cdot A\cdot U_1U_2\cdots U_r = I
\end{equation*}
Di sini $I$ adalah matriks identitas parsial,
$D_0$ adalah matriks diagonal,
$L_i$ adalah matriks kombinasi baris segitiga bawah,
dan $U_j$ adalah matriks kombinasi kolom segitiga atas.

Semua operasi baris dapat dibatalkan
(sebagai contoh, $2\rho_1+\rho_2$
dibatalkan dengan $-2\rho_1+\rho_2$).
Dengan demikian, setiap matriks segitiga bawah tersebut
memiliki invers.
Demikian pula, setiap matriks segitiga atas memiliki invers.
Artinya, kita dapat melakukan manipulasi aljabar untuk mengisolasi~$A$
dalam persamaan itu.
Oleh karena itu,
jika reduksi Gauss--Jordan pada matriks~$A$ tidak memerlukan pertukaran,
kita memperoleh faktorisasi matriks awal
\begin{equation*} 
  A = L_1^{-1}\cdots L_s^{-1}\cdot D\cdot U_r^{-1}\cdots U_1^{-1}
\end{equation*}
dengan $D$ merupakan hasil kali $D_0^{-1}$ dan identitas parsial~$I$,
yang jelas merupakan matriks diagonal.

Untuk menangani masalah pertukaran, kita dapat melakukan pertukaran terlebih dahulu:~sebelum
memfaktorkan matriks awal, pertukarkan dahulu baris-barisnya
dengan matriks permutasi~$P$.
\begin{equation*}
  P\cdot A = L_1^{-1}\cdots L_s^{-1}\cdot D\cdot U_r^{-1}\cdots U_1^{-1}
  \tag{$*{*}*$}
\end{equation*}

Sekarang tibalah langkah penutup.
Invers matriks segitiga bawah merupakan matriks segitiga bawah,
dan hasil kali matriks-matriks segitiga bawah juga
merupakan matriks segitiga bawah. Karena itu, semua $L_i^{-1}$ dalam~($*{*}*$)
dapat digabungkan menjadi satu matriks segitiga bawah~$L$.
Demikian pula, semua $U_i^{-1}$ dapat digabungkan menjadi satu
matriks segitiga atas~$U$.
Jadi, sebarang matriks memiliki faktorisasi $PA=LDU$.

Kita mengilustrasikannya dengan transformasi umum $\nbyn{2}$ pada $\Re^2$ yang,
terhadap basis standar, direpresentasikan sebagai berikut.
\begin{equation*}
  T=
  \begin{mat}
    1 &2 \\
    3 &4
  \end{mat}
\end{equation*}
Penerapan Metode Gauss di sini cukup sederhana.
\begin{equation*}
  \begin{mat}
    1 &2 \\
    3 &4
  \end{mat}
  \grstep{-3\rho_1+\rho_2}  
  \begin{mat}
    1 &2 \\
    0 &-2
  \end{mat}
  \grstep{-(1/2)\rho_2}  
  \begin{mat}
    1 &2 \\
    0 &1
  \end{mat}
  \grstep{-2\chi_1+\chi_2}  
  \begin{mat}
    1 &0 \\
    0 &1
  \end{mat}
\end{equation*}
(Kita menuliskan $\chi_i$ untuk kolom.)
Berikut adalah faktorisasi yang bersesuaian.
\begin{equation*}
  \begin{mat}
    1 &2 \\
    3 &4
  \end{mat}
  =
  \begin{mat}
   1 &0 \\
   3 &1 
  \end{mat}
  \begin{mat}
    1 &0 \\
    0 &-2
  \end{mat}
  \begin{mat}
    1 &2 \\
    0 &1
  \end{mat}
\end{equation*}
Hasil kalinya sesuai.
\begin{sagecommandline}
sage: L = matrix(QQ, [[1, 0], [3, 1]])
sage: D = matrix(QQ, [[1, 0], [0, -2]])
sage: U = matrix(QQ, [[1, 2], [0, 1]])
sage: L*D*U  
\end{sagecommandline}

Kita membahas ini untuk memahami
efek geometris transformasi umum tersebut.
\begin{equation*}
  \begin{mat}
   1 &2 \\
   3 &4 
  \end{mat}
  \colvec{x \\ y}
\end{equation*}
\begin{sagecommandline}
sage: load("plot_action.sage")
sage: q = plot_square_action(1,0,0,1) 
sage: q.set_axes_range(-2, 4, -1, 7) 
sage: q.save("graphics/geo300a.pdf")
sage: p = plot_square_action(1,2,3,4) 
sage: p.set_axes_range(-2, 4, -1, 7) 
sage: p.save("graphics/geo300b.pdf")
\end{sagecommandline}
\begin{sagesilent}
load("plot_action.sage")
q = plot_square_action(1,0,0,1) 
q.set_axes_range(-2, 4, -1, 7) 
q.save("graphics/geo300a.pdf")
p = plot_square_action(1,2,3,4) 
p.set_axes_range(-2, 4, -1, 7) 
p.save("graphics/geo300b.pdf")
\end{sagesilent}
\begin{equation*}
  \vcenteredhbox{\includegraphics{graphics/geo300a.pdf}}
  \quad\mapsvia{\big (\begin{smallmatrix} 1 &2 \\ 3 &4 \end{smallmatrix}\big )}\quad
  \vcenteredhbox{\includegraphics{graphics/geo300b.pdf}}
\end{equation*}
\noindent Uraikan matriks itu dengan faktorisasi~$LDU$ di atas.
\begin{equation*}
  \begin{mat}
    1 &2 \\
    3 &4
  \end{mat}
  \colvec{x \\ y}
  =
  \begin{mat}
   1 &0 \\
   3 &1 
  \end{mat}
  \begin{mat}
    1 &0 \\
    0 &-2
  \end{mat}
  \begin{mat}
    1 &2 \\
    0 &1
  \end{mat}  
  \colvec{x \\ y}
\end{equation*}
Pertama-tama, kita akan melihat efek $U$, $D$, dan~$L$ secara terpisah.
Kemudian kita akan melihat efek kumulatifnya:~aksi~$U$, lalu~$DU$,
dan akhirnya~$LDU$.

Matriks yang diterapkan pertama, yaitu matriks paling kanan~$U$,
merupakan geseran sejajar sumbu~$x$.
\begin{sagecommandline}
sage: load("plot_action.sage")
sage: q = plot_square_action(1,0,0,1) 
sage: q.set_axes_range(-2, 4, -2.25, 4.25) 
sage: q.save("graphics/geo303a.pdf")
sage: p = plot_square_action(1,2,0,1) 
sage: p.set_axes_range(-2, 4, -2.25, 4.25) 
sage: p.save("graphics/geo303b.pdf")
\end{sagecommandline}
\begin{sagesilent}
load("plot_action.sage")
q = plot_square_action(1,0,0,1) 
q.set_axes_range(-2, 4, -2.25, 4.25) 
q.save("graphics/geo303a.pdf")
p = plot_square_action(1,2,0,1) 
p.set_axes_range(-2, 4, -2.25, 4.25) 
p.save("graphics/geo303b.pdf")
\end{sagesilent}
\begin{equation*}
  \vcenteredhbox{\includegraphics{graphics/geo303a.pdf}}
  \quad\mapsvia{\big (\begin{smallmatrix} 1 &2 \\ 0 &1 \end{smallmatrix}\big )}\quad
  \vcenteredhbox{\includegraphics{graphics/geo303b.pdf}}
\end{equation*}
Matriks kedua~$D$ mengubah skala dan orientasi.
\begin{sagecommandline}
sage: load("plot_action.sage")
sage: q = plot_square_action(1,0,0,1) 
sage: q.set_axes_range(-2, 4, -2.25, 4.25) 
sage: q.save("graphics/geo302a.pdf")
sage: p = plot_square_action(1,0,0,-2) 
sage: p.set_axes_range(-2, 4, -2.25, 4.25) 
sage: p.save("graphics/geo302b.pdf")
\end{sagecommandline}
\begin{sagesilent}
load("plot_action.sage")
q = plot_square_action(1,0,0,1) 
q.set_axes_range(-2, 4, -2.25, 4.25) 
q.save("graphics/geo302a.pdf")
p = plot_square_action(1,0,0,-2) 
p.set_axes_range(-2, 4, -2.25, 4.25) 
p.save("graphics/geo302b.pdf")
\end{sagesilent}
\begin{equation*}
  \vcenteredhbox{\includegraphics{graphics/geo302a.pdf}}
  \quad\mapsvia{\big (\begin{smallmatrix} 1 &0 \\ 0 &-2 \end{smallmatrix}\big )}\quad
  \vcenteredhbox{\includegraphics{graphics/geo302b.pdf}}
\end{equation*}

Matriks paling kiri, $L$, merupakan geseran sejajar sumbu~$y$.
\begin{sagecommandline}
sage: q = plot_square_action(1,0,0,1) 
sage: q.set_axes_range(-2, 4, -2.25, 4.25) 
sage: q.save("graphics/geo301a.pdf")
sage: p = plot_square_action(1,0,3,1) 
sage: p.set_axes_range(-2, 4, -2.25, 4.25) 
sage: p.save("graphics/geo301b.pdf")
\end{sagecommandline}
\begin{sagesilent}
load("plot_action.sage")
q = plot_square_action(1,0,0,1) 
q.set_axes_range(-2, 4, -2.25, 4.25) 
q.save("graphics/geo301a.pdf")
p = plot_square_action(1,0,3,1) 
p.set_axes_range(-2, 4, -2.25, 4.25) 
p.save("graphics/geo301b.pdf")
\end{sagesilent}
\begin{equation*}
  \vcenteredhbox{\includegraphics{graphics/geo301a.pdf}}
  \quad\mapsvia{\big (\begin{smallmatrix} 1 &0 \\ 3 &1 \end{smallmatrix}\big )}\quad
  \vcenteredhbox{\includegraphics{graphics/geo301b.pdf}}
\end{equation*}
  
Matriks segitiga bawah dan segitiga atas tidak
mengubah orientasi.
Setiap perubahan orientasi dalam $LDU$ terjadi melalui entri diagonal yang
negatif.
(Perhatikan bahwa jika suatu matriks memerlukan pertukaran baris, sehingga $PA=LDU$, situasinya
lebih rumit.
Setiap pertukaran baris membalik orientasi, dari berlawanan arah jarum jam menjadi searah jarum jam, atau
sebaliknya.
Jadi,
jika matriks permutasi memerlukan jumlah pertukaran yang ganjil,
matriks itu mengubah orientasi,
sedangkan jumlah pertukaran yang genap mempertahankan orientasi.)

Kita mengakhiri pembahasan $LDU$ dengan memperhatikan efek kumulatifnya.
Perhatikan secara berurutan aksi $U$, lalu~$DU$, dan kemudian~$LDU$.
Sebagai contoh, berikut adalah efek kumulatif pemetaan-pemetaan tersebut pada
sudut kanan atas persegi satuan.
\begin{equation*}
  \colvec{1 \\ 1}
  \mapsunder{U}
  \colvec{3 \\ 1}
  \mapsunder{D}
  \colvec{3 \\ -2}
  \mapsunder{L}
  \colvec{3 \\ 7}
\end{equation*}
Berikut adalah gambar untuk seluruh persegi.
Sekali lagi kita melihat geseran sejajar sumbu~$x$, diikuti perubahan skala dan
orientasi, lalu geseran sejajar sumbu~$y$.
\begin{sagecommandline}
sage: L = matrix(QQ, [[1, 0], [3, 1]])
sage: D = matrix(QQ, [[1, 0], [0, -2]])
sage: U = matrix(QQ, [[1, 2], [0, 1]])
sage: DU = D*U
sage: LDU = L*DU  
sage: load("plot_action.sage")
sage: p = plot_square_action(1,0,0,1) 
sage: p.set_axes_range(-4, 4, -2, 7) 
sage: p.save("graphics/geo304a.pdf")
sage: p = plot_square_action(U[0][0], U[0][1], U[1][0], U[1][1]) 
sage: p.set_axes_range(-4, 4, -2, 7) 
sage: p.save("graphics/geo304b.pdf")
sage: p = plot_square_action(DU[0][0], DU[0][1], DU[1][0], DU[1][1]) 
sage: p.set_axes_range(-4, 4, -2, 7) 
sage: p.save("graphics/geo304c.pdf")
sage: p = plot_square_action(LDU[0][0], LDU[0][1], LDU[1][0],LDU[1][1]) 
sage: p.set_axes_range(-4, 4, -2, 7) 
sage: p.save("graphics/geo304d.pdf")
\end{sagecommandline}
\begin{sagesilent}
L = matrix(QQ, [[1, 0], [3, 1]])
D = matrix(QQ, [[1, 0], [0, -2]])
U = matrix(QQ, [[1, 2], [0, 1]])
DU = D*U
LDU = L*DU  
load("plot_action.sage")
p = plot_square_action(1,0,0,1) 
p.set_axes_range(-4, 4, -2, 7) 
p.save("graphics/geo304a.pdf")
p = plot_square_action(U[0][0], U[0][1], U[1][0], U[1][1]) 
p.set_axes_range(-4, 4, -2, 7) 
p.save("graphics/geo304b.pdf")
p = plot_square_action(DU[0][0], DU[0][1], DU[1][0], DU[1][1]) 
p.set_axes_range(-4, 4, -2, 7) 
p.save("graphics/geo304c.pdf")
p = plot_square_action(LDU[0][0], LDU[0][1], LDU[1][0], LDU[1][1]) 
p.set_axes_range(-4, 4, -2, 7) 
p.save("graphics/geo304d.pdf")
\end{sagesilent}
\begin{center}
  \begin{tabular}{rcl}
    \vcenteredhbox{\includegraphics{graphics/geo304a.pdf}}
    &$\mapsvia{\big (\begin{smallmatrix} 1 &2 \\ 0 &1 \end{smallmatrix}\big )}$
    &\vcenteredhbox{\includegraphics{graphics/geo304b.pdf}}  \\
    &$\mapsvia{\big (\begin{smallmatrix} 1 &0 \\ 0 &-2 \end{smallmatrix}\big )}$
    &\vcenteredhbox{\includegraphics{graphics/geo304c.pdf}}  \\
    &$\mapsvia{\big (\begin{smallmatrix} 1 &0 \\ 3 &1 \end{smallmatrix}\big )}$
    &\vcenteredhbox{\includegraphics{graphics/geo304d.pdf}} 
  \end{tabular} 
\end{center}



\section{Kode sumber plot\_action.sage}
Rutin \inlinecode{plot_square_action}
menerima empat entri matriks $\nbyn{2}$
dan mengembalikan sebuah daftar grafik.
\lstinputlisting[firstline=92,lastline=103]{plot_action.sage}

Sebagian besar pekerjaan
dilakukan oleh rutin pembantu \inlinecode{color_square_list}.
\lstinputlisting[firstline=53,lastline=90]{plot_action.sage}
Ada dua hal teknis yang perlu
dijelaskan.
\inlinecode{ZORDER} menentukan urutan \Sage{} dalam
menggambar objek; di sini kita ingin persegi satuan digambar
setelah sumbu agar warna-warnanya terlihat.
Selain itu, sambungan antarruas garis tampak kurang rapi,
sehingga sambungan itu kita tutupi dengan sebuah titik.

\endinput


Perlu dikerjakan:
